%% ---------------------------------------------------------------------
\section{Discussion}\label{sec:discussion}

\paragraph{What the pumping lemma buys}
\Cref{lem:pumping} turns an infinite family of requirements on the
reflected triangles into constraints on the half-line diagram only. It is
satisfied by Mazoyer's solution, as it must be, and it eliminates, at
negligible cost, the whole class of rules whose half-line becomes periodic
behind the front too early. \Cref{thm:third} shows that the constraint has a
transparent geometric meaning: in every minimal-time solution the
non-periodic part of the half-line reaches the line $i=t/3$. For Mazoyer's
solution every depth-row is $3$-periodic from time $2j+1$ or $2j+2$, i.e.\
the periodic zone behind the front is bounded by a signal of speed $1/2$,
comfortably within the bound.

\paragraph{Why lucky rules are hard to kill}
Small sets of necessary conditions are typically met by rules that are not
solutions; the question is how many lengths are needed before all such rules
disappear (if no solution exists, some finite set of lengths suffices, since
there are finitely many rules). For four states the answer is $9$ (or $8$ with a long
half-line). For five states the rule space is larger by roughly $38$
orders of magnitude ($5^{94}$ versus $4^{45}$ transition tables), and our
experiments show partial rules for all lengths up to at least $12$, and for
any single additional length we tried. Under the pumping constraints the
surviving rules are no longer random: their half-lines exhibit a speed-$1/3$
boundary, and the remaining freedom lies in the reflected triangles.

\paragraph{The mirrored construction}
The cone of \cref{lem:cone} ends where the anti-diagonal $t+i=2n-2$ meets
the line $i=t/3$. By \cref{lem:determinacy}, the part of $R_n$ whose cone lies
in the periodic zone behind the front is the same for all large $n$ of a
given residue class, and the information that makes the right end fire
exactly at time $2n-2$ enters $R_n$ near that crossing and travels to the
right end at maximal speed. Heuristically, the right half of every
reflected triangle therefore has to behave like a \emph{mirrored}
minimal-time synchronization, started by the reflection and running on the
background left behind the front instead of on quiescent cells. A
five-state solution would have to realize both constructions---the direct
one on the quiescent background and the mirrored one on the front's
background---with the same four working states. We regard this tension as
the structural reason why five states are hard; turning it into a theorem
(for instance a mirrored version of \cref{thm:third}, for which the missing
ingredient is a border at the far end of the mirrored problem) is the most
promising route to a conceptual proof.

\paragraph{Open problems}
\begin{enumerate}[label=(\arabic*)]
\item Does a five-state minimal-time solution exist? By
  \cref{thm:frontier}, a negative answer along the lines of \cref{thm:four}
  requires lines of length at least $13$.
\item Is there a necessary condition, in the spirit of \cref{lem:pumping},
  that constrains the reflected triangles themselves and not only the
  half-line?
\item What is the smallest $N$ such that some five-state rule is
  $\{2,\dots,N\}$-minimal but no five-state rule is
  $\{2,\dots,N+1\}$-minimal, if such an $N$ exists?
\end{enumerate}
